\newpage
\section{Kompaktheit}
\label{2.3}
\begin{definition}{Kompaktheit}
    Eine Menge $K\subset X$ heißt kompakt, wenn jede offene Überdeckung von K eine endliche Teilüberdeckung besitzt.
\end{definition}

\begin{satz}{Heine-Borel}
    Eine Menge \(K\subset\mathbb{R}^n\) ist genau dann kompakt, wenn sie abgeschlossen und beschränkt ist.
\end{satz}

\begin{satz}{Folgenkompaktheit}
    Eine Menge $K\subset X$ ist kompakt genau dann, wenn jede Folge in $K$ eine konvergente Teilfolge besitzt, deren Grenzwert in $K$ liegt.
\end{satz}


\begin{satz}{Kompakt $\Rightarrow$ abgeschlossen und beschränkt}
    In einem metrischen Raum ist jede kompakte Menge abgeschlossen und beschränkt.
\end{satz}

\begin{bemerkungbox}
    Die Umkehrung von Satz 2.28 ist im Allgemeinen falsch. Im diskreten metrischen Raum $(\mathbb{R}^n,d)$ ist jede Menge abgeschlossen und beschränkt, aber $\mathbb{N}$ ist nicht kompakt.
\end{bemerkungbox}

\begin{satz}{Heine--Borel}
    Sei $(\mathbb{R}^n,d)$ mit der euklidischen Metrik. Dann sind äquivalent:
    \begin{itemize}
        \item[(i)] $K$ ist kompakt,Korollar 2.33
        \item[(ii)] $K$ ist abgeschlossen und beschränkt.
    \end{itemize}
\end{satz}

\begin{satz}{Hilfsatz}
    Sei $Q=\{x\in\mathbb{R}^n : a_i\le x_i\le b_i,\ i=1,\dots,n\}$ ein abgeschlossener Quader. Dann ist $Q$ kompakt.
\end{satz}

\begin{satz}{Überdeckungskompakt $\Leftrightarrow$ folgenkompakt}
    Eine Menge $K$ in einem metrischen Raum ist genau dann kompakt, wenn jede Folge in $K$ eine konvergente Teilfolge besitzt, deren Grenzwert in $K$ liegt.
\end{satz}

\begin{satz}{}
    Sei $(X,d)$ metrischer Raum, $K\subset X$ kompakt und $A\subset K$ abgeschlossen. Dann ist $A$ kompakt.
\end{satz}